\documentclass[11pt]{article}
\usepackage[margin=1in]{geometry}
\usepackage{amsmath,amssymb,amsthm}
\usepackage{graphicx}
\usepackage{array}
\usepackage{xr}
\externaldocument[P-]{tit-stationary}
\usepackage[hidelinks]{hyperref}
\def\UrlBreaks{\do\_\do\.\do\-\do\/}

\newcommand{\Th}[1]{Theorem~\ref{P-#1}}
\newcommand{\Pp}[1]{Proposition~\ref{P-#1}}
\newcommand{\Co}[1]{Corollary~\ref{P-#1}}
\newcommand{\Se}[1]{Section~\ref{P-#1}}

\title{What One Description Conveys to Observers of a Signal.\\[4pt]
\large An Observation Theory Companion to ``Rate, Leakage, and Distortion for Stationary Gaussian Sources''}
\author{Andrew H. Bond\\ San Jos\'e State University}
\date{Draft, November 2026}

\begin{document}
\maketitle

\begin{abstract}
Observation Theory treats every party that holds a description as an
observer with its own view of the source, its own geometry for weighing
errors, and its own budget, and it asks what the same description
conveys to each. A first companion read a paper on memoryless Gaussian
sources in that language. This one reads its sequel on stationary
sources. The sequel's theorems are stated and proved there in standard
information theory, and nothing here adds to them. What time adds to
the picture is that an observer of a signal has a view at every
frequency and a history it may or may not keep. Four statements follow.
Two observers of a stationary signal meet frequency by frequency and
share one price. They can read disjoint bands of the same signal. A
delay in an observer's view conveys nothing by itself. And for a
description of one moment of a scalar Gaussian context, everything an
observer has kept enters through one number. The last two say that staleness is a matter of what
an observer holds when it uses its view, whether it kept less or the
rest has not yet arrived, and that a delay alone costs nothing.
\end{abstract}

\section{Observers of a signal}
\label{sec:intro}

In Observation Theory an observer is a triple $O=(C,G,B)$: what it reads
of the source, the geometry in which it weighs errors, and its budget
\cite{obsrep}. The first companion \cite{supp1} applied this to one
description of a memoryless Gaussian vector held by two observers, the
decoder it was written for and a holder it was not written for. The
content conveyed to the holder is the conditional mutual information
between the source and the description given the holder's view, which
secure source coding calls leakage \cite{villard2013}.

A signal adds time. The source is a stationary process, the decoder
reconstructs part of it sample by sample, and the holder's view is a
filtered and noisy copy of the process. Two things are new for an
observer. Its view $C$ is now a filter, so it has a view at every
frequency. And it has a history, of which it may keep all, a recent
part, or one sample. The companion paper \cite{paper2} works out the
rate and the leakage per symbol in this setting. The table records the
correspondence.

\begin{center}
\begin{tabular}{>{\raggedright\arraybackslash}p{7.2cm}>{\raggedright\arraybackslash}p{7.2cm}}
\hline
Observation Theory & Companion paper \\
\hline
observer the description is written for & decoder, reading $Y_t=FT_t$ under weight $W$ \\
observer the description is not written for & holder, seeing $S=h\ast T+U$ \\
the second observer's view at a frequency & $H(\omega)$ and $\Sigma_U(\omega)$ \\
the second observer's geometry at a frequency & $J(\omega)=H(\omega)^{\mathsf H}\Sigma_U(\omega)^{-1}H(\omega)$ \\
content conveyed per symbol & leakage per symbol \\
what the second observer keeps & the index set $\mathcal A$ of retained samples \\
\hline
\end{tabular}
\end{center}

Every entry on the right is standard, and renaming does not create a
result. As in the first companion, the contribution of Observation
Theory is that it asked the questions, here by treating the holder as an
observer with a view over time, and that it gives the answers one
reading.

\section{Two observers meet frequency by frequency and share one price}
\label{sec:price}

\Th{thm:periodic} and \Th{thm:spectral} say that the frontier of rate
and conveyed content is solved one frequency at a time. At each
frequency the optimal description is reverse water-filling on the
decoder's weight tilted by what the holder's view adds, exactly as for
a memoryless source, and the frequencies are tied together by one
number, the water level.

Read as a statement about observers, the two geometries interact only
within a frequency. What the holder sees at one frequency has no effect
on how the description should treat another, except through the single
price that the decoder's budget sets for distortion. The first
companion found that the region splits along subspaces the two
geometries share. Stationarity gives every pair of observers such
subspaces without asking anything of their geometries. Time-invariance
of the source and of both views is what the two observers have in
common, and it is enough.

\section{Two observers can read disjoint bands of one signal}
\label{sec:bands}

For a scalar signal with a scalar context, \Th{thm:pair} gives the order
in which frequencies are described. The description written for the
decoder alone starts where the signal's spectrum is largest. The
description that conveys least to the holder starts where that spectrum
times one plus the holder's coherent signal-to-noise ratio is largest.
\Co{cor:bands} and the instance of \Se{sec:examples} show that the two
can describe disjoint bands at the same budget. There the description
written for the decoder occupies the low band and the one that conveys
least to the holder occupies the high band, and each costs the other
party a factor of $2.6$ in conveyed content or $1.7$ in rate.

This is the frequency-domain form of the first companion's central
finding, that the second observer's geometry changes what an optimal
description reads. For memoryless sources the change was a direction
that turns with the budget. For a signal it can be a different band.
The reason is the same identity. The content conveyed to the holder is
the rate minus what the description tells the holder about its own
view. A band the holder already sees well is cheap to describe, however
little of the signal's power it carries.

\section{A delay conveys nothing by itself}
\label{sec:delay}

Suppose the holder's view of the context arrives late. One might expect
a stale view to help the holder less, so that the same description
conveys more to it. \Pp{prop:allpass} says that it does not. If the
holder has its whole observation, a delay of any length, and more
generally any invertible filtering of what it sees, leaves the region
unchanged, for circular blocks of every length and in the spectral
limit. An ordinary finite window of a stationary signal has edges, and
there the statement holds only up to the samples a delay moves across
them.

The reason is plain once stated. An observer that has everything loses
nothing when what it holds is passed through an invertible map. What
the result rules out is a reading in which the age of a view is itself
a cost. It does not say that timing never matters. An observer that
must use its view before all of it has arrived holds less, and that is
the next section.

\section{For one moment, what an observer has kept is one number}
\label{sec:kept}

The cost appears when the observer keeps less. \Pp{prop:retention}
treats a description of one moment of the signal, with a holder that
keeps any set of its past and future samples. The context there is a
scalar Gaussian process. Everything the holder has kept
enters through the conditional variance $q$ of the context at that
moment given the kept samples, and the content conveyed is the scalar
closed form of the memoryless theory at an effective noise $q/(1-q)$.
The one number is a feature of the scalar context. For a vector context
the same argument would leave a conditional covariance matrix in its
place.

Three cases show the range. An observer that keeps one sample of age
$\Delta$ has a $q$ that rises with $\Delta$, and the content conveyed to
it rises with the age in closed form. An observer that keeps its whole
past, and is allowed to wait $\Delta$ steps beyond the moment
described, has a $q$ that falls geometrically in $\Delta$ from the
filtering variance to the smoothing variance. An observer that keeps
everything has the smoothing variance, and $\Delta$ does not appear.
For the instance of \Se{sec:examples} the three values of $q$ and the
conveyed content are listed there.

Together with Section~\ref{sec:delay}, this says that staleness is a
matter of what an observer holds at the time it uses its view. It can
hold less because it kept less, as with one sample, or because the rest
has not arrived, as with a causal prefix. An observer that kept its
entire past still lacks the future, and the fixed-lag case measures
exactly that. The observer triple has a view, a
geometry, and a budget. A signal makes the view a function of what is
retained, and the one-moment result shows that, for Gaussian signals,
retention acts on the content conveyed only by rescaling the observer's
noise.

\section{What does not need Observation Theory}
\label{sec:standard}

The method of the companion paper is classical. Splitting a circulant
problem over Fourier frequencies and passing to a spectral limit is how
the rate--distortion function of a Gaussian process is derived, and the
limit over Toeplitz windows is Szeg\H{o}'s theorem. The optimization
at each frequency rests on the memoryless theorems of \cite{paper1}.
The theorems for the stationary signal itself need more than that: an
order structure and a subadditivity argument for finite windows, the
approximation of the signal by one of finite memory, a coupling at the
encoder, guard intervals between coded blocks, the degradation of the
holder by added noise, control of the observations at the edges of a
window, and the removal of the encoder's randomness. All of it is
information theory and none of it is Observation Theory. The
convergence statements (\Th{thm:spectral},
\Th{thm:windows}), the window converse (\Pp{prop:window}), and the
coding theorem (\Th{thm:achieve}) would be stated the same way
by a reader who had never heard of observers. The fixed-lag formula in
\Pp{prop:retention} is the smoothing recursion of estimation theory.

\section{Open questions about observers over time}
\label{sec:open}

\emph{An observer with a causal view of a block.} The one-number result
is for a description of one moment. For a block description against a
holder that keeps only a causal prefix, the per-symbol content is not
known. This is the natural place for a theory of observers that cannot
wait.

\emph{An observer written for under a delay constraint.} The decoder in
the companion paper may wait for the whole block. A decoder that must
reconstruct with bounded delay is a different first observer, and the
frontier against a holder is open.

\emph{An observer whose view has unbounded memory.} When the holder
sees the signal through a filter of finite length, in noise with a
continuous positive definite spectral density, the spectral region is
the operational region of the stationary signal for every continuous
positive definite spectral density of the signal, by a converse
through windows and a coding theorem with deterministic encoders
(\Th{thm:achieve}, \Th{thm:filter}). In that statement the holder has
its observations at the times of the coded window. If it has its whole
observation process, before and after the window as well, the coding
theorem still holds and the matching converse is open. This is a
different matter from the one-moment result of
Section~\ref{sec:kept}, where the holder may keep any of its past and
future samples and the answer is complete. When the holder's filter has an
infinite impulse response, the region is the limit of periodic sources,
and its status for the signal itself is open.

\emph{Several observers.} With several holders the weighted objective
has one tilt per holder at each frequency, by the same argument. Whether
the resulting surface is achievable, and whether three observers can
read three disjoint bands, is open.

\section*{Where this fits}

The observer triple and the rate--content region are defined in
\cite{obsrep}. The memoryless two-observer theory is \cite{paper1}, read
in \cite{supp1}. The scripts and the notebook that produce the numbers
quoted here are listed in the companion paper's appendix on numerical
checks.

\begin{thebibliography}{9}
\bibitem{paper2} A. H. Bond, ``Rate, leakage, and distortion for
stationary Gaussian sources: Water-filling on a tilted spectrum,''
draft, 2026.
\bibitem{paper1} A. H. Bond, ``Rate, leakage, and distortion for Gaussian
sources: Water-filling on a tilted weight,'' submitted to \emph{IEEE
Trans. Inf. Theory}, 2026, archived at
\url{https://doi.org/10.5281/zenodo.23127802}.
\bibitem{supp1} A. H. Bond, ``What one description conveys to each
observer: An Observation Theory companion to `Rate, leakage, and
distortion for Gaussian sources','' draft, 2026,
\texttt{paper/tit-rate-leakage/supplement-ot.tex} in
\texttt{https://github.com/ahb-sjsu/geometric-observation}.
\bibitem{obsrep} A. H. Bond, ``Observer representation,'' draft, 2026,
\texttt{paper/observer-representation.tex} in the same repository.
\bibitem{villard2013} J. Villard and P. Piantanida, ``Secure multiterminal
source coding with side information at the eavesdropper,'' \emph{IEEE
Trans. Inf. Theory}, vol. 59, no. 6, pp. 3668--3692, 2013.
\end{thebibliography}

\end{document}
